\documentclass[10pt]{extarticle}
\usepackage[a4paper,margin=1.5cm]{geometry}
\usepackage[T1]{fontenc}
\usepackage{amsmath,amssymb}
\usepackage{graphicx}
\usepackage[export]{adjustbox}
\usepackage{multicol}
\usepackage{enumitem}
\usepackage{booktabs,array,tabularx}
\setlength{\columnsep}{0.7cm}
\setlength{\parindent}{0pt}
\setlength{\parskip}{3pt}
\pagestyle{empty}
\makeatletter
% figure registry (figdefs.tex is generated by the extract step)
\newcommand{\qfigdef}[4]{\@namedef{qf@file@#1}{#2}\@namedef{qf@w@#1}{#3}\@namedef{qf@wide@#1}{#4}}
\newcommand{\qf@check}[1]{\@ifundefined{qf@file@#1}{\PackageError{score2}{Unknown figure id #1}{}}{}}
\newcommand{\qf@img}[4][]{\includegraphics[width=#2\dimexpr\@nameuse{qf@w@#4}pt\relax,max width=#3,max height=0.45\textheight,#1]{\@nameuse{qf@file@#4}}}
% Full-width questions: finalize wraps every question that contains a wide figure in \qfullbegin...\qfullend,
% so the WHOLE question (stem, figures, options) is set across the page in reading order, then the two columns
% resume. The question's head - its number, stem and first figure (\qfullsplit marks the end of the head) -
% is kept on one page, so a number is never left at a page foot with its figure overleaf; the rest may break
% across pages normally, so no big gaps are left. Without a \qfullsplit, the whole question is kept together
% if it is under 45% of a page. Never write these yourself.
\newif\ifqf@full
\newif\ifqf@boxed
\newsavebox{\qf@box}
\newcommand{\qf@open}{\global\qf@fulltrue\par\global\qf@boxedtrue
  \setbox\qf@box\vbox\bgroup\hsize\textwidth\linewidth\textwidth}
\newcommand{\qf@emit}[1]{\par\egroup\global\qf@boxedfalse
  \ifdim\dimexpr\ht\qf@box+\dp\qf@box\relax<#1\textheight\noindent\usebox{\qf@box}\par\else\unvbox\qf@box\fi}
\newcommand{\qfullsplit}{\ifqf@boxed\qf@emit{0.85}\fi}
\newcommand{\qf@close}{\ifqf@boxed\qf@emit{0.45}\fi\global\qf@fullfalse}
\newcommand{\qfullbegin}{\end{multicols}\qf@open}
\newcommand{\qfullend}{\qf@close\begin{multicols}{2}}
\newcommand{\qfullstart}{\qf@open}
\newcommand{\qfulldone}{\qf@close}
% \qfig[scale]{id}: block figure, centred, directly where it is written (size set by the tool).
% Fallback if a wide figure is not inside a full-width question: break out of the columns in place - no floats.
\newcommand{\qfig}[2][1]{\qf@check{#2}%
  \ifqf@full\qf@blk{#1}{#2}\else
  \ifnum\@nameuse{qf@wide@#2}=1 \expandafter\@firstoftwo\else\expandafter\@secondoftwo\fi
  {\end{multicols}\begin{center}\qf@img{#1}{\textwidth}{#2}\end{center}\begin{multicols}{2}}%
  {\qf@blk{#1}{#2}}\fi}
\newcommand{\qf@blk}[2]{\par\smallskip{\centering\qf@img{#1}{\linewidth}{#2}\par}\smallskip}
% \qfigi[scale]{id}: inline figure (inside option lists / tables / minipages); never breaks columns.
\newcommand{\qfigi}[2][1]{\qf@check{#2}\qf@img[valign=c]{#1}{\linewidth}{#2}}
\makeatother
% structure macros - lint reads these to check completeness and figure placement
\newcommand{\papersec}[1]{\par\bigskip{\centering\Large\bfseries PAPER-#1\par}\medskip}
\newcommand{\partsec}[2]{\par\medskip\noindent{\large\bfseries PART-#1: #2}\par\smallskip}
\newcommand{\sect}[1]{\par\smallskip\noindent\textbf{#1}\par}
\newcommand{\qn}[1]{\par\smallskip\noindent\textbf{#1.}~}
\newcommand{\paraq}[2]{\par\smallskip\noindent\textbf{Paragraph for Questions #1 and #2}\par}
\newenvironment{qoptions}{\begin{itemize}[leftmargin=1.8em,itemsep=1pt,topsep=2pt,label={}]}{\end{itemize}}
\newcommand{\opt}[2]{\item[(#1)] #2}
% figure registry (generated by score2_tools.py finalize)
\qfigdef{p1_phy_s1_q02_1}{Score Advanced - Test 40 (27 Apr 2025)-img/p1_phy_s1_q02_1.png}{139.6}{0}
\qfigdef{p1_phy_s1_q03_1}{Score Advanced - Test 40 (27 Apr 2025)-img/p1_phy_s1_q03_1.png}{80.1}{0}
\qfigdef{p1_phy_s2_para3_4_1}{Score Advanced - Test 40 (27 Apr 2025)-img/p1_phy_s2_para3_4_1.png}{66.2}{0}
\qfigdef{p1_phy_s3_q07_1}{Score Advanced - Test 40 (27 Apr 2025)-img/p1_phy_s3_q07_1.png}{95.3}{0}
\qfigdef{p1_phy_s3_q07_2}{Score Advanced - Test 40 (27 Apr 2025)-img/p1_phy_s3_q07_2.png}{51.9}{0}
\qfigdef{p1_phy_s3_q07_3}{Score Advanced - Test 40 (27 Apr 2025)-img/p1_phy_s3_q07_3.png}{67.3}{0}
\qfigdef{p1_phy_s3_q07_4}{Score Advanced - Test 40 (27 Apr 2025)-img/p1_phy_s3_q07_4.png}{51.4}{0}
\qfigdef{p1_phy_s3_q07_5}{Score Advanced - Test 40 (27 Apr 2025)-img/p1_phy_s3_q07_5.png}{52.9}{0}
\qfigdef{p1_phy_s3_q08_1}{Score Advanced - Test 40 (27 Apr 2025)-img/p1_phy_s3_q08_1.png}{172.2}{0}
\qfigdef{p1_phy_s3_q08_2}{Score Advanced - Test 40 (27 Apr 2025)-img/p1_phy_s3_q08_2.png}{84.7}{0}
\qfigdef{p1_phy_s3_q08_3}{Score Advanced - Test 40 (27 Apr 2025)-img/p1_phy_s3_q08_3.png}{148.8}{0}
\qfigdef{p1_phy_s3_q08_4}{Score Advanced - Test 40 (27 Apr 2025)-img/p1_phy_s3_q08_4.png}{98.1}{0}
\qfigdef{p1_phy_s4_q01_1}{Score Advanced - Test 40 (27 Apr 2025)-img/p1_phy_s4_q01_1.png}{135.4}{0}
\qfigdef{p1_phy_s4_q05_1}{Score Advanced - Test 40 (27 Apr 2025)-img/p1_phy_s4_q05_1.png}{101.8}{0}
\qfigdef{p1_phy_s4_q06_1}{Score Advanced - Test 40 (27 Apr 2025)-img/p1_phy_s4_q06_1.png}{175.7}{0}
\qfigdef{p1_chem_s1_q01_1}{Score Advanced - Test 40 (27 Apr 2025)-img/p1_chem_s1_q01_1.png}{130.6}{0}
\qfigdef{p1_chem_s2_para1_2_1}{Score Advanced - Test 40 (27 Apr 2025)-img/p1_chem_s2_para1_2_1.png}{220.6}{0}
\qfigdef{p1_chem_s2_para3_4_1}{Score Advanced - Test 40 (27 Apr 2025)-img/p1_chem_s2_para3_4_1.png}{186.3}{0}
\qfigdef{p1_chem_s3_q07_1}{Score Advanced - Test 40 (27 Apr 2025)-img/p1_chem_s3_q07_1.png}{80.8}{0}
\qfigdef{p1_chem_s3_q07_2}{Score Advanced - Test 40 (27 Apr 2025)-img/p1_chem_s3_q07_2.png}{83.7}{0}
\qfigdef{p1_chem_s3_q07_3}{Score Advanced - Test 40 (27 Apr 2025)-img/p1_chem_s3_q07_3.png}{79.4}{0}
\qfigdef{p1_chem_s4_q01_1}{Score Advanced - Test 40 (27 Apr 2025)-img/p1_chem_s4_q01_1.png}{51.2}{0}
\qfigdef{p1_chem_s4_q01_3}{Score Advanced - Test 40 (27 Apr 2025)-img/p1_chem_s4_q01_3.png}{54.0}{0}
\qfigdef{p1_chem_s4_q01_2}{Score Advanced - Test 40 (27 Apr 2025)-img/p1_chem_s4_q01_2.png}{58.4}{0}
\qfigdef{p1_chem_s4_q01_4}{Score Advanced - Test 40 (27 Apr 2025)-img/p1_chem_s4_q01_4.png}{49.0}{0}
\qfigdef{p1_chem_s4_q02_2}{Score Advanced - Test 40 (27 Apr 2025)-img/p1_chem_s4_q02_2.png}{59.0}{0}
\qfigdef{p1_chem_s4_q02_3}{Score Advanced - Test 40 (27 Apr 2025)-img/p1_chem_s4_q02_3.png}{50.8}{0}
\qfigdef{p1_chem_s4_q02_1}{Score Advanced - Test 40 (27 Apr 2025)-img/p1_chem_s4_q02_1.png}{63.0}{0}
\qfigdef{p1_chem_s4_q02_4}{Score Advanced - Test 40 (27 Apr 2025)-img/p1_chem_s4_q02_4.png}{52.9}{0}
\qfigdef{p1_chem_s4_q02_5}{Score Advanced - Test 40 (27 Apr 2025)-img/p1_chem_s4_q02_5.png}{64.4}{0}
\qfigdef{p1_chem_s4_q02_7}{Score Advanced - Test 40 (27 Apr 2025)-img/p1_chem_s4_q02_7.png}{65.8}{0}
\qfigdef{p1_chem_s4_q02_8}{Score Advanced - Test 40 (27 Apr 2025)-img/p1_chem_s4_q02_8.png}{65.1}{0}
\qfigdef{p1_chem_s4_q02_6}{Score Advanced - Test 40 (27 Apr 2025)-img/p1_chem_s4_q02_6.png}{63.0}{0}
\qfigdef{p2_phy_s1_q02_1}{Score Advanced - Test 40 (27 Apr 2025)-img/p2_phy_s1_q02_1.png}{51.0}{0}
\qfigdef{p2_phy_s1_q03_1}{Score Advanced - Test 40 (27 Apr 2025)-img/p2_phy_s1_q03_1.png}{84.3}{0}
\qfigdef{p2_phy_s1_q04_1}{Score Advanced - Test 40 (27 Apr 2025)-img/p2_phy_s1_q04_1.png}{202.1}{0}
\qfigdef{p2_phy_s2_para1_2_1}{Score Advanced - Test 40 (27 Apr 2025)-img/p2_phy_s2_para1_2_1.png}{103.0}{0}
\qfigdef{p2_phy_s2_para3_4_1}{Score Advanced - Test 40 (27 Apr 2025)-img/p2_phy_s2_para3_4_1.png}{130.6}{0}
\qfigdef{p2_phy_s2_para5_6_1}{Score Advanced - Test 40 (27 Apr 2025)-img/p2_phy_s2_para5_6_1.png}{127.2}{0}
\qfigdef{p2_phy_s2_q05_1}{Score Advanced - Test 40 (27 Apr 2025)-img/p2_phy_s2_q05_1.png}{133.1}{0}
\qfigdef{p2_phy_s2_q06_1}{Score Advanced - Test 40 (27 Apr 2025)-img/p2_phy_s2_q06_1.png}{136.9}{0}
\qfigdef{p2_phy_s3_q07_1}{Score Advanced - Test 40 (27 Apr 2025)-img/p2_phy_s3_q07_1.png}{118.6}{0}
\qfigdef{p2_phy_s3_q08_1}{Score Advanced - Test 40 (27 Apr 2025)-img/p2_phy_s3_q08_1.png}{153.0}{0}
\qfigdef{p2_phy_s4_q02_1}{Score Advanced - Test 40 (27 Apr 2025)-img/p2_phy_s4_q02_1.png}{109.9}{0}
\qfigdef{p2_phy_s4_q03_1}{Score Advanced - Test 40 (27 Apr 2025)-img/p2_phy_s4_q03_1.png}{87.9}{0}
\qfigdef{p2_phy_s4_q04_1}{Score Advanced - Test 40 (27 Apr 2025)-img/p2_phy_s4_q04_1.png}{149.5}{0}
\qfigdef{p2_phy_s4_q06_1}{Score Advanced - Test 40 (27 Apr 2025)-img/p2_phy_s4_q06_1.png}{133.2}{0}
\qfigdef{p2_chem_s1_q01_1}{Score Advanced - Test 40 (27 Apr 2025)-img/p2_chem_s1_q01_1.png}{135.4}{0}
\qfigdef{p2_chem_s1_q03_1}{Score Advanced - Test 40 (27 Apr 2025)-img/p2_chem_s1_q03_1.png}{202.2}{0}
\qfigdef{p2_chem_s3_q07_1}{Score Advanced - Test 40 (27 Apr 2025)-img/p2_chem_s3_q07_1.png}{93.9}{0}
\qfigdef{p2_chem_s3_q07_5}{Score Advanced - Test 40 (27 Apr 2025)-img/p2_chem_s3_q07_5.png}{37.2}{0}
\qfigdef{p2_chem_s3_q07_2}{Score Advanced - Test 40 (27 Apr 2025)-img/p2_chem_s3_q07_2.png}{112.1}{0}
\qfigdef{p2_chem_s3_q07_6}{Score Advanced - Test 40 (27 Apr 2025)-img/p2_chem_s3_q07_6.png}{55.4}{0}
\qfigdef{p2_chem_s3_q07_3}{Score Advanced - Test 40 (27 Apr 2025)-img/p2_chem_s3_q07_3.png}{99.2}{0}
\qfigdef{p2_chem_s3_q07_7}{Score Advanced - Test 40 (27 Apr 2025)-img/p2_chem_s3_q07_7.png}{30.1}{0}
\qfigdef{p2_chem_s3_q07_4}{Score Advanced - Test 40 (27 Apr 2025)-img/p2_chem_s3_q07_4.png}{96.2}{0}
\qfigdef{p2_chem_s3_q07_8}{Score Advanced - Test 40 (27 Apr 2025)-img/p2_chem_s3_q07_8.png}{33.5}{0}
\qfigdef{p2_chem_s4_q06_1}{Score Advanced - Test 40 (27 Apr 2025)-img/p2_chem_s4_q06_1.png}{176.6}{0}
\raggedcolumns
\makeatletter\@beginparpenalty=10000\makeatother
\newcolumntype{L}[1]{>{\raggedright\arraybackslash}p{#1}}
\begin{document}
\begin{center}{\Large\bfseries Score Advanced --- Test 40 (27 Apr 2025)}\end{center}
\begin{multicols}{2}
\raggedcolumns
\papersec{1}
\partsec{1}{PHYSICS}
\sect{SECTION-I}
\qn{1} The vessel volume is V, the initial pressure is P$_0$. The process is assumed to be isothermal, and the evacuation rate is constant \& equal to C and independent of pressure. [The evacuation rate is the gas volume being evacuated per unit time with that volume being measured under the gas pressure attained by that moment]. Choose the \textbf{CORRECT} option(s).
\begin{qoptions}
\opt{A}{Pressure P of gas at time `t' will be P$_0$e$^{-\frac{\mathrm{Ct}}{\mathrm V}}$}
\opt{B}{Pressure P of gas at time `t' will be P$_0-\dfrac{\mathrm{Ct}}{\mathrm V}$}
\opt{C}{No. of moles evacuated in time `t' will be linearly proportional to time t.}
\opt{D}{No. of moles evacuated in time `t' will be exponentially dependent on time t.}
\end{qoptions}
\qn{2} A glass partially filled with water is fastened to a wedge that slides, without friction, down a large plane inclined at an angle $\alpha$ (as shown in the figure). The mass of the inclined plane is M, the combined mass of the wedge, the glass and the water is m. If there were no motion the water surface would be horizontal. Choose the \textbf{CORRECT} option(s).\par
\qfig{p1_phy_s1_q02_1}
\begin{qoptions}
\opt{A}{If the inclined plane is fixed, then the angle free surface of water will make with the inclined plane in steady state will be zero.}
\opt{B}{If the inclined plane is fixed, then the angle surface of water will make with the inclined plane in steady state will depend on m.}
\opt{C}{If the inclined plane can move freely in the horizontal direction, then the angle surface of water will make with the inclined plane in steady state will be zero.}
\opt{D}{If the inclined plane can move freely in the horizontal direction, then the angle surface of water will make with the inclined plane in steady state will depend on m \& M.}
\end{qoptions}
\qn{3} In an old record found in a laboratory, a process 1$\to$2$\to$3 was shown. Over the time, ink faded and it became impossible to see the pressure and volume axes. However, descriptions given there reveal that the states 1 and 3 lie on an isochore corresponding to a volume V, amount of heat supplied during the whole process 1$\to$2$\to$3 is zero and the gas involved is one mole of helium, then choose the \textbf{CORRECT} option(s).\par
\qfig{p1_phy_s1_q03_1}
\begin{qoptions}
\opt{A}{Volume occupied by the gas in the state 2 will depend on pressure at 2.}
\opt{B}{Volume occupied by the gas in the state 2 is independent of pressure at 2.}
\opt{C}{Volume occupied by the gas in the state 2 is 4V.}
\opt{D}{Heat supplied will be positive if we complete cycle 1$\to$2$\to$3$\to$1.}
\end{qoptions}
\qn{4} A point mass of 1 kg moving at constant speed of 5 m/s on an elliptical path experiences a centripetal force of 20N when at either end point of the major axis and a similar force of 2.5 N at each end of minor axis. Then :-
\begin{qoptions}
\opt{A}{length of semi major axis is 5 m}
\opt{B}{length of semi major axis is 10 m}
\opt{C}{length of semi minor axis is 2.5 m}
\opt{D}{length of semi minor axis is 5 m}
\end{qoptions}
\sect{SECTION-II(i)}
\paraq{1}{2}
There are three sound sources producing sound with slight different frequency. The equation of sound waves when received by detector is given by\par
$\Delta\mathrm P_1=\Delta\mathrm P_0\sin298\pi\left(\mathrm t-\dfrac{\mathrm x}{\mathrm v}\right)$\par
$\Delta\mathrm P_2=\Delta\mathrm P_0\sin300\pi\left(\mathrm t-\dfrac{\mathrm x}{\mathrm v}\right)$\par
$\Delta\mathrm P_3=\Delta\mathrm P_0\sin302\pi\left(\mathrm t-\dfrac{\mathrm x}{\mathrm v}\right)$\par
Where $\Delta$P is change in pressure and v is speed of sound in medium. Let the detector is placed at x = 0 and intensity due to each source is I$_0$ at the place of detector.
\qn{1} Number of maxima local as well as global of loudness in every one second will be :
\qn{2} If maximum intensity of resulting sound is 9I$_0$ and detector cannot record an intensity lower than 4I$_0$, If the duration of time for which detector remains idle during first one second is $\tau$(sec), find 3$\tau$.
\paraq{3}{4}
A flat conducting sheet A is suspended by an insulating thread at distance `s' from both the surfaces formed by the bend conducting sheet B as shown in the figure. The entire space between bent sheet B is filled with medium of dielectric constant K. The sheets A and B are oppositely charged, the difference in potential is constant and equal to $\Delta$V. This causes a force F, pulling A downward. (Neglect gravity)\par
\qfig{p1_phy_s2_para3_4_1}
\qn{3} If the modulus of work done by external agent needed to slowly increase the inserted distance by $\Delta$y is $\dfrac{\mathrm P(\Delta\mathrm V)^2\mathrm{bK}\varepsilon_0\Delta\mathrm y}{\mathrm s}$. Find the value of P.
\qn{4} An expression for the difference in potential $\Delta$V in terms of F and relevant dimensions shown in the figure is $\sqrt{\dfrac{\mathrm{PsF}}{\varepsilon_0\mathrm{bK}}}$. Find the value of P.
\paraq{5}{6}
Hydrogen atoms in ground state are excited by means of a monochromatic radiation of wavelength 969$\cdot$69 \AA. After absorbing the energy of radiation, hydrogen atoms go to the excited state. After 10$^{-8}$ s, the hydrogen atom will come to the ground state by emitting the absorbed energy. Different wavelengths are possible in the spectrum. (Take : R = 1.1 $\times$ 10$^7$ / m)
\qn{5} The electron of excited hydrogen atom is in which of the energy state n?
\qn{6} Number of different wavelengths present in the spectrum is :-
\sect{SECTION-II(ii)}
\qn{7} A very small circular object is kept infront of an optical device as shown in figure. The plane of object is parallel to the optical device. One quadrant of circle is hatched as shown. Match the images as seen by the observer (ignoring magnification).\par
\qfig{p1_phy_s3_q07_1}
{\footnotesize\setlength{\tabcolsep}{1.5pt}\noindent\begin{tabular}{|c|L{4.0cm}|c|c|}
\hline
 & \multicolumn{1}{>{\centering\arraybackslash}p{4.0cm}|}{List-I (Device)} & & \multicolumn{1}{c|}{List-II (Probable Image)} \\ \hline
(P) & Plane mirror (The observer is at A) & (1) & \qfigi{p1_phy_s3_q07_2} \\ \hline
(Q) & Concave mirror (The observer is at A) object between F and 2F & (2) & \qfigi{p1_phy_s3_q07_3} \\ \hline
(R) & Convex mirror (The observer is at A) & (3) & \qfigi{p1_phy_s3_q07_4} \\ \hline
(S) & Convex lens (The observer is at B) object between F and optical centre & (4) & \qfigi{p1_phy_s3_q07_5} \\ \hline
\end{tabular}}\par
\qn{8} Match statement of List-I with List-II and select the correct answer using the code given below the list.\par
{\footnotesize\setlength{\tabcolsep}{1.5pt}\noindent\begin{tabular}{|c|L{5.2cm}|c|L{1cm}|}
\hline
 & \multicolumn{1}{>{\centering\arraybackslash}p{5.2cm}|}{List-I} & & \multicolumn{1}{c|}{List-II} \\ \hline
(P) & \qfigi{p1_phy_s3_q08_1}\par Graph-1 and Graph-2 represents the displacement time graph of two identical particles, performing SHM with same total energy. The value of a$_1$/a$_2$ is \underline{\hspace{0.8cm}}. & (1) & 0 \\ \hline
(Q) & \qfigi{p1_phy_s3_q08_2}\par A circuit contains two capacitors as shown in the figure. Initially switch S$_1$ is closed and S$_2$ is open. At t = 0, S$_2$ is closed keeping S$_1$ closed, and a charge Q$_1$ is supplied by the battery to the circuit. After long time, S$_1$ is opened keeping S$_2$ closed. Now battery supplied charge Q$_2$. The value of Q$_2$/Q$_1$ is & (2) & 1 \\ \hline
(R) & \qfigi{p1_phy_s3_q08_3}\par A massive rope of uniform mass density and length $\ell$ is suspended from points A and B in vertical plane as shown in the figure. Point C is the lowest point of the chord. Given length of rope from A to C = $\ell_1$ and length of rope from B to C = $\ell_2$. The value of $\ell_2/\ell_1$ is & (3) & 2 \\ \hline
(S) & \qfigi{p1_phy_s3_q08_4}\par A particle performs SHM along a straight line between the points M and N about O(mean position) as shown in figure. t$_1$ and t$_2$ are the time taken by the particle from M to P and P to N respectively. Find t$_1$/t$_2$ (given OP = PN) & (4) & 3 \\ \hline
\end{tabular}}\par
\sect{SECTION-III}
\qn{1} Parallel conducting rails are located horizontally at a distance L from each other and placed in a uniform magnetic field whose induction vector B is directed at an angle of 37$^\circ$ to their plane as shown. The rails are connected by a resistor of high resistance R. Far from the resistor on the rails lies massive conducting rod of mass m and length L. It is perpendicular to the rails. What force (in N) do you need to exert on the rod in the horizontal direction so that it slides along the rails with constant speed v? The surface is rough with coefficient of friction $\mu$ = 0.5 and B = 0.5 T, L = 30 cm, m = 1 kg, R = 27 m$\Omega$, v = 10 m/s, g = 10 m/s$^2$.\par
\qfig{p1_phy_s4_q01_1}
\qn{2} An ac voltmeter with large impedance is connected one by one across the inductor, the capacitor, and the resistor in a series circuit having an alternating emf of 125 V (rms); the meter gives the same reading in volts in each case. The angular frequency of the AC source is 4$\sqrt3$ rad/s. Now the circuit is disconnected from the AC source, the capacitor is separately charged and reconnected to the inductor and the resistor, all three elements in series. The circuit is found to undergo underdamped oscillations. What is the angular frequency of oscillation in rad/s ? Round off to nearest integer if necessary.
\qn{3} $^4$He with energy E$_0$ bombard stationary $^7$Li nuclei. They become composite nuclei $^{11}$B after completely inelastic collision. It further splits into $^{10}$B and neutron $^1$n. The reaction equation is $^4$He + $^7$Li $\to$ ($^{11}$B) $\to$ $^{10}$B + $^1$n. Assume that the energy E$_0$ is minimum. What is the corresponding neutron kinetic energy E$_\mathrm n$ (in MeV)? Mass of helium = 4.0026 u, Mass of Lithium = 7.0160 u, Mass of $^{11}$B = 11.0093 u, mass of neutron = 1.0087 u, Mass of $^{10}$B = 10.0129 u, 1 amu = 930 MeV. Find 50 E$_\mathrm n$. Round off to nearest integer if required.
\qn{4} Intensity observed in an interference pattern is I = I$_0$ sin$^2$ $\theta$. At $\theta$ = 30$^\circ$, intensity I = 5 $\pm$ 0.0020 W/m$^2$. If I$_0$ = 20 W/m$^2$ and percentage error in angle $\theta$ is $\dfrac{\mathrm n}{\pi}\sqrt3\times10^{-2}\%$, then the value of n is:
\qn{5} A simple pendulum of inextensible string of length $\ell$ and mass m is hinged at point A in a vertical plane as shown in the figure. The ball is released from the same horizontal level as that of A and at a distance $\dfrac{\sqrt3\ell}2$ from point `A'. At the lowest point bob hits a block of identical mass m elastically. Find the maximum compression in the spring ($\dfrac{\mathrm m\ell}{\mathrm k}=\dfrac{56}{10}$kg$-$m$^2$/N and ground is smooth). Height of A from ground is equal to $\ell$.\par
\qfig{p1_phy_s4_q05_1}
\qn{6} In the given circuit, the value of R (in $\Omega$) so that thermal power generated in R will be maximum.\par
\qfig{p1_phy_s4_q06_1}
\raggedcolumns
\papersec{1}
\partsec{2}{CHEMISTRY}
\sect{SECTION-I}
\qn{1} \qfig{p1_chem_s1_q01_1}
Select correct statements for all above alkenes.
\begin{qoptions}
\opt{A}{Compound `e' shows the highest heat of combustion.}
\opt{B}{Compound `d' is the thermodynamically least reactive towards catalytic hydrogenation.}
\opt{C}{On ozonolysis using O$_3$/Me$_2$S, only 2 alkenes give smallest non$-$cyclic ketone.}
\opt{D}{Only 2 alkenes among these give chiral alcohol product(s) on reactive with B$_2$H$_6$/THF followed by H$_2$O$_2$/OH$^-$.}
\end{qoptions}
\qn{2} How many of these solutions (of equal volume) on mixing create a Buffer solution?
\begin{qoptions}
\opt{A}{0.1 M NaOH, 0.3 M HCN}
\opt{B}{0.3 M H$_2$SO$_4$, 0.8 M CH$_3$COO$^-$Na$^+$}
\opt{C}{0.2 M Ba(OH)$_2$, 0.4 M HCN}
\opt{D}{0.3 M H$_2$CO$_3$, 0.1 M NaOH}
\end{qoptions}
\qn{3} Among the following solutions, which two solutions will show `highest' and `lowest' increase in temperature on mixing?\par
Assume that initially all solutions at 25$^\circ$C and density of solutions are identical.
\begin{qoptions}
\opt{A}{0.4 M, 50 mL NaOH with 0.5 M, 100 mL HNO$_3$}
\opt{B}{0.2 M, 75 mL H$_2$SO$_4$ with 0.6 M, 25 mL NaOH}
\opt{C}{0.3 M, 20 mL HBr with 0.3 M, 80 mL Ba(OH)$_2$}
\opt{D}{0.9 M, 40 mL Ba(OH)$_2$ with 0.7 M, 60 mL HI}
\end{qoptions}
\qn{4} Select correct for metallurgical process which is/are involve in extraction metal from galena.
\begin{qoptions}
\opt{A}{Reduction of roasted ore with fresh galena}
\opt{B}{Copper impurities can be remove by liquation}
\opt{C}{Alkali metal cynide selectively prevent galena coming to the froth}
\opt{D}{Oxidation of silver impurity in molten metal by O$_2$}
\end{qoptions}
\sect{SECTION-II(i)}
\paraq{1}{2}
The following amino acids are used to form dipeptides :\par
\qfig{p1_chem_s2_para1_2_1}
\qn{1} A dipeptide `X' is formed which has maximum possible degree of unsaturation. The number of $\pi-$bonds in that dipeptide is 2n. Find the value of n.
\qn{2} Another dipeptide `Y' with only one Sulphur atom and two nitrogen atoms is formed. Let `P' is the number of oxygen atoms in `Y' and `Q' is the number of chiral centres in `Y'. Then report the value of (P + Q)
\paraq{3}{4}
\qfig{p1_chem_s2_para3_4_1}
Consider the following graph of `Bond order' vs `total number of protons' in diatomic species `a' to `f' According to molecular orbital theory, if bond order is zero at `a', `c' and `g' (protons in `g' = 20') then species is highly unstable \& most likely non$-$existent.
\qn{3} Among the following species, how many will either relate to `d' or `e'?\par
Ne$_2^+$, NO$^+$, N$_2^+$, N$_2$, O$_2$, O$_2^{2+}$, O$_2^-$, CN$^-$, BN
\qn{4} If X$_2$ diatomic species relates to `b' and Y$_2$ relates to `f', then total number of electrons in $\sigma$*2s of both X$_2$ \& Y$_2$ species for ground state only.
\paraq{5}{6}
16 g of Sulphur is reacted with `X' g of O$_2$ to produce both SO$_2$ \& SO$_3$ gases. The actual proportion of SO$_2$ \& SO$_3$ produced depends on value of `X' (relative to mass of Sulphur taken). It is best to associate combustion of Sulphur to SO$_2$ \& SO$_3$ using two separate chemical equations. (S = 32, O = 16)
\qn{5} Find mass \% of SO$_2$ in product mixture if X = 25 (refer initial data in comprehension) and no sulphur is left behind.
\qn{6} If value of X = 20, then the ratio of moles of SO$_2$ \& SO$_3$ formed as products, if none of the reactant is left behind is a : 1. The value of `a' is.
\sect{SECTION-II(ii)}
\qn{7} Correct matching from the two lists.\par
{\footnotesize\setlength{\tabcolsep}{1.5pt}\noindent\begin{tabular}{|c|L{4.2cm}|c|L{2.4cm}|}
\hline
 & \multicolumn{1}{>{\centering\arraybackslash}p{4.2cm}|}{List-I} & & \multicolumn{1}{>{\centering\arraybackslash}p{2.4cm}|}{List-II} \\ \hline
(P) & \qfigi{p1_chem_s3_q07_1} & (1) & Positional isomers \\ \hline
(Q) & CH$_3$CH$_2$CH$_2$OH \& CH$_3$OCH$_2$CH$_3$ & (2) & Metamers \\ \hline
(R) & \qfigi{p1_chem_s3_q07_2} & (3) & Chain isomers \\ \hline
(S) & \qfigi{p1_chem_s3_q07_3} & (4) & Functional isomers \\ \hline
\end{tabular}}\par
\qn{8} Correct matching from the two lists.\par
{\footnotesize\setlength{\tabcolsep}{1.5pt}\noindent\begin{tabular}{|c|L{1.6cm}|c|L{4.8cm}|}
\hline
 & \multicolumn{1}{>{\centering\arraybackslash}p{1.6cm}|}{List-I} & & \multicolumn{1}{>{\centering\arraybackslash}p{4.8cm}|}{List-II (Type of $\pi-$bonds)} \\ \hline
(P) & NO$_3^-$ & (1) & Only p$\pi-$d$\pi$ bonding \\ \hline
(Q) & (CH$_3$)$_2$O & (2) & No p$\pi-$d$\pi$ and p$\pi-$p$\pi$ bonding \\ \hline
(R) & SO$_3$ & (3) & Only p$\pi-$p$\pi$ bonding \\ \hline
(S) & SO$_4^{2-}$ & (4) & Both p$\pi-$p$\pi$ and p$\pi-$d$\pi$ bonding \\ \hline
\end{tabular}}\par
\sect{SECTION-III}
\qn{1} How many of these do not give self$-$aldol reaction?\par
(i) CH$_3$CHO\quad (ii) CH$_3$COCH$_2$CH$_3$\quad (iii) CD$_3$COCD$_3$\quad (iv) \qfigi{p1_chem_s4_q01_1}\quad (v) Ph$-$CHO\quad (vi) HCHO\quad \qfigi{p1_chem_s4_q01_3}\quad \qfigi{p1_chem_s4_q01_2}\quad \qfigi{p1_chem_s4_q01_4}
\qn{2} How many of these can exhibit geometrical isomerism?\par
\qfigi{p1_chem_s4_q02_2}\quad \qfigi{p1_chem_s4_q02_3}\par \qfigi{p1_chem_s4_q02_1}\quad \qfigi{p1_chem_s4_q02_4}\quad \qfigi{p1_chem_s4_q02_5}\par (vi) (CH$_3$)$_2$C=CHMe\quad \qfigi{p1_chem_s4_q02_7}\par \qfigi{p1_chem_s4_q02_8}\quad \qfigi{p1_chem_s4_q02_6}
\qn{3} A weak acid HA aqueous solution is found to have conductivity value of 4.5 $\times$ 10$^{-4}$ ohm$^{-1}$cm$^{-1}$. Its concentration was found to be 0.01 molar and conductivity of water is found to be 3 $\times$ 10$^{-4}$ Scm$^{-1}$. If K$_\mathrm a$ of this acid HA is y $\times$ 10$^{-6}$, report value of `y'? (Answer to the nearest integer)\par
\textbf{Given:} $\Lambda^0_{\mathrm{m(HA)}}$ = 500 Scm$^2$mol$^{-1}$ [Limiting molar conductance]
\qn{4} A sample containing fixed amount of ideal gas mass, is subjected to 50\% increase in T \& 25\% decrease in pressure. If the ratio of initial density over final density of gas sample is x : 1, the value of `x' is?
\qn{5} If reaction A $\xrightarrow{\mathrm K}$ B has K = 6.93 $\times$ 10$^{-3}$ min$^{-1}$ at 20$^\circ$C, then half-life of `A' at 70$^\circ$C is $\dfrac{25\mathrm x}{16}$ min. Find value of `x', if temperature coefficient of reaction is assumed to have a constant value of 2 ?
\qn{6} SO$_3$ decomposes into SO$_2$ \& O$_2$ according to this equation in a rigid closed vessel of volume 1 litre.\par
SO$_3$(g)$\longrightarrow$SO$_2$(g) + $\dfrac12$O$_2$(g)\par
At time = 20 min, rate of appearance of O$_2$ was 20 g L$^{-1}$ per minute and $\dfrac{\mathrm d[\mathrm{SO}_3]}{\mathrm{dt}}$ at that instant is reported as $\dfrac{-\mathrm x}3$g L$^{-1}$s$^{-1}$. Then value of `x' is? [Atomic mass of S = 32]
\raggedcolumns
\papersec{1}
\partsec{3}{MATHEMATICS}
\sect{SECTION-I}
\qn{1} Which of the following is/are correct
\begin{qoptions}
\opt{A}{$^{30}\mathrm C_0\cdot{}^{90}\mathrm C_{60}-{}^{30}\mathrm C_1\cdot{}^{89}\mathrm C_{60}+{}^{30}\mathrm C_2\,{}^{88}\mathrm C_{60}+\ldots+{}^{30}\mathrm C_{30}\,{}^{60}\mathrm C_{60}={}^{60}\mathrm C_{30}$}
\opt{B}{$^{30}\mathrm C_0\cdot{}^{90}\mathrm C_{30}-{}^{30}\mathrm C_1\cdot{}^{89}\mathrm C_{30}+{}^{30}\mathrm C_2\,{}^{88}\mathrm C_{30}+\ldots+{}^{30}\mathrm C_{30}\,{}^{60}\mathrm C_{30}=0$}
\opt{C}{$^{30}\mathrm C_0\cdot{}^{90}\mathrm C_{30}-{}^{30}\mathrm C_1\cdot{}^{88}\mathrm C_{30}+{}^{30}\mathrm C_2\,{}^{86}\mathrm C_{30}+\ldots+{}^{30}\mathrm C_{30}\,{}^{30}\mathrm C_{30}=(2)^{30}$}
\opt{D}{$^{30}\mathrm C_0\cdot{}^{90}\mathrm C_{30}-{}^{30}\mathrm C_1\cdot{}^{87}\mathrm C_{30}+{}^{30}\mathrm C_2\,{}^{84}\mathrm C_{30}+\ldots+{}^{30}\mathrm C_{20}\,{}^{30}\mathrm C_{30}=1$}
\end{qoptions}
\qn{2} Let $-$x + y + 2 = 0 and 2x $-$ y + 6 = 0 be tangent of parabola having focus at (1, 3) then :
\begin{qoptions}
\opt{A}{Equation of directrix of parabola is 3x + 4y = 11}
\opt{B}{Length of latus rectum of parabola is $\left(\dfrac45\right)$}
\opt{C}{Equation of tangent at vertex is 3x + 4y $-$ 13 = 0}
\opt{D}{Equation of axis of parabola is 4x $-$ 3y + 5 = 0}
\end{qoptions}
\qn{3} Which of the following is (are) correct (where [.] represent GIF and \{.\} represent fractional part function)
\begin{qoptions}
\opt{A}{Let f(x) = $\left\{\dfrac{\mathrm x}3\right\}+\left[\dfrac{\mathrm x}5\right]$, x $\in$ (0, 301) then number of points where f(x) is discontinuous is 120}
\opt{B}{Let f(x) = [3 sin x + 4 cos x], x $\in$ [0, 2$\pi$] then number of points where f(x) is discontinuous is 18}
\opt{C}{$\displaystyle\int_0^\pi|3\sin\mathrm x+4\cos\mathrm x|\,\mathrm{dx}$ = 10}
\opt{D}{Let f(x) = $\dfrac1{\mathrm x-1}+\dfrac1{\mathrm x-2}+\dfrac1{\mathrm x-3}+\dfrac1{\mathrm x-4}$ then number of real roots of equation f(x) = 0 is 3}
\end{qoptions}
\qn{4} If A$_\mathrm i$ is the area bounded by $|$x $-$ a$_\mathrm i|$ + $|$y$|$ = b$_\mathrm i$, i $\in$ N where a$_{\mathrm i+1}$ = a$_\mathrm i$ + $\dfrac32$b$_\mathrm i$ and b$_{\mathrm i+1}$ = $\dfrac{\mathrm b_\mathrm i}2$, a$_1$ = 0, b$_1$= 32 then
\begin{qoptions}
\opt{A}{A$_3$ = 128}
\opt{B}{A$_3$ = 256}
\opt{C}{$\displaystyle\lim_{\mathrm n\to\infty}\sum_{\mathrm i=1}^{\mathrm n}\mathrm A_\mathrm i=\dfrac83(32)^2$}
\opt{D}{$\dfrac{\mathrm A_6}{\mathrm A_{10}}$ = 256}
\end{qoptions}
\sect{SECTION-II(i)}
\paraq{1}{2}
Let x$^2-$ x $-$ 1 = 0 is a factor of ax$^9$ + bx$^8$ + 1 = 0 then answer the following questions :
\qn{1} The value of `2a+b' is :
\qn{2} Number of real roots of equation ax$^9$ + bx$^8$ + 1 = 0 is:
\paraq{3}{4}
Let `z$_1$' and `z$_2$' be any two distinct complex number which satisfy z$^{120}$ = 1.
\qn{3} The probability that $|$z$_1-$ z$_2|$ $\geqslant$ $\dfrac{\sqrt5+1}2$ is $\dfrac{\mathrm p}{\mathrm q}$ (where p and q are co-prime) then the value of q $-$ 2p is :
\qn{4} The number of possible ordered pair (z$_1$, z$_2$) such that $|$z$_1-$ z$_2|$ $\leqslant$ $\sqrt{2-\sqrt3}$
\paraq{5}{6}
Consider all the permutation of the word `MAHAKUMBH'
\qn{5} Number of permutation of the word in which vowels are in alphabetical order is P then the value of $\dfrac{\mathrm P}{30}$ is :
\qn{6} Number of permutation of the word if no two alike letters are together is N then total number of divisor of N is :
\sect{SECTION-II(ii)}
\qn{7} Let A be any 3 $\times$ 3 matrix such that $|$A$|$ = 2, then match the column\par
{\footnotesize\setlength{\tabcolsep}{1.5pt}\noindent\begin{tabular}{|c|L{3.6cm}|c|L{2.4cm}|}
\hline
 & \multicolumn{1}{>{\centering\arraybackslash}p{3.6cm}|}{List-I} & & \multicolumn{1}{>{\centering\arraybackslash}p{2.4cm}|}{List-II} \\ \hline
(P) & $|3\,\mathrm{adj}(-2\mathrm{adj}(\mathrm{adj}(\mathrm A)))|$ & (1) & 2$^{14}\cdot$3$^3$ \\ \hline
(Q) & $|2\,\mathrm{adj}(3\mathrm{adj}(\mathrm{adj}(2\mathrm A)))|$ & (2) & 3$^{12}\cdot$2$^8$ \\ \hline
(R) & $|\mathrm{adj}(\mathrm{adj}(\mathrm{adj}(3\mathrm A)))|$ & (3) & 2$^{35}\cdot$3$^6$ \\ \hline
(S) & $|\mathrm{adj}(\mathrm{adj}(3\mathrm{adj}(\mathrm A)))|$ & (4) & 3$^{24}\cdot$2$^8$ \\ \hline
\end{tabular}}\par
\qn{8} Match the following\par
{\footnotesize\setlength{\tabcolsep}{1.5pt}\noindent\begin{tabular}{|c|L{5.0cm}|c|L{0.8cm}|}
\hline
 & \multicolumn{1}{>{\centering\arraybackslash}p{5.0cm}|}{List-I} & & \multicolumn{1}{c|}{List-II} \\ \hline
(P) & If $|\vec{\mathrm a}|=|\vec{\mathrm b}|=|\vec{\mathrm c}|$, angle between each pair of vectors is $\dfrac\pi3$ and $|\vec{\mathrm a}+\vec{\mathrm b}+\vec{\mathrm c}|=\sqrt6$, then $2|\vec{\mathrm a}|$ is equal to & (1) & 3 \\ \hline
(Q) & If $\vec{\mathrm a}$ is perpendicular to $\vec{\mathrm b}+\vec{\mathrm c}$, $\vec{\mathrm b}$ is perpendicular to $\vec{\mathrm c}+\vec{\mathrm a}$, $\vec{\mathrm c}$ is perpendicular $\vec{\mathrm a}+\vec{\mathrm b}$, where $|\vec{\mathrm a}|$ = 2, $|\vec{\mathrm b}|$ = 3 and $|\vec{\mathrm c}|$ = 6 then $|\vec{\mathrm a}+\vec{\mathrm b}+\vec{\mathrm c}|-2$ is equal to & (2) & 2 \\ \hline
(R) & $\vec{\mathrm a}=2\hat{\mathrm i}+3\hat{\mathrm j}-\hat{\mathrm k}$, $\vec{\mathrm b}=-\hat{\mathrm i}+2\hat{\mathrm j}-4\hat{\mathrm k}$, $\vec{\mathrm c}=\hat{\mathrm i}+\hat{\mathrm j}+\hat{\mathrm k}$ and $\vec{\mathrm d}=3\hat{\mathrm i}+2\hat{\mathrm j}+\hat{\mathrm k}$, then $\dfrac17(\vec{\mathrm a}\times\vec{\mathrm b})\cdot(\vec{\mathrm c}\times\vec{\mathrm d})$ is equal to & (3) & 4 \\ \hline
(S) & If $|\vec{\mathrm a}|=|\vec{\mathrm b}|=|\vec{\mathrm c}|=2$ and $\vec{\mathrm a}\cdot\vec{\mathrm b}=\vec{\mathrm b}\cdot\vec{\mathrm c}=\vec{\mathrm c}\cdot\vec{\mathrm a}=2$ then $[\vec{\mathrm a}\,\vec{\mathrm b}\,\vec{\mathrm c}]\cos45^\circ$ is equal to & (4) & 5 \\ \hline
\end{tabular}}\par
\sect{SECTION-III}
\qn{1} If $\displaystyle\lim_{\mathrm n\to\infty}\sum_{\mathrm n=1}^{\mathrm n}\tan^{-1}\left(\dfrac{2\mathrm n}{\mathrm n^4+\mathrm n^2+2}\right)$ = L , then the value of [8L] ( where [ . ] is represent G.I.F )
\qn{2} Let f(x) = $\sqrt{7-3\mathrm x}$, x $\leqslant$ $\dfrac73$ then number of solutions of the equation f(x) = f$^{-1}$(x) is :
\qn{3} Number of solution of the equation log$_\mathrm k$x = cos($\pi$x) is 3 , then sum of all possible integers value of `k' is :
\qn{4} The value of integral $\displaystyle\int_0^1\dfrac{12\mathrm x^{11}}{3\mathrm x^6-3\mathrm x^3+1}$dx is :
\qn{5} Let P = (1)$^{1/4}$ + (2)$^{1/4}$ +\ldots.+(625)$^{1/4}$ and $\left[\dfrac{\mathrm P}5\right]$ = N then total number of even divisor of N is (where [.] represents greatest integer function).
\qn{6} The value of $\displaystyle\int_0^2\left(\sqrt{1+\mathrm x^3}+\sqrt[3]{\mathrm x^2+2\mathrm x}\right)$dx is equal to :
\raggedcolumns
\papersec{2}
\partsec{1}{PHYSICS}
\sect{SECTION-I}
\qn{1} The expression for the angular momentum of a point mass about origin is given by $\vec{\mathrm L}$ = 2$\alpha$t$\,\ell$n$\left(\dfrac{\mathrm t}{\mathrm t+1}\right)\hat{\mathrm k}$ where $\alpha$ is a certain constant and `t' is time in seconds.
\begin{qoptions}
\opt{A}{The torque acting on the mass as a function of time is 2$\alpha\left[\dfrac1{\mathrm t+1}+\ell\mathrm n\left(\dfrac{\mathrm t}{\mathrm t+1}\right)\right]\hat{\mathrm k}$.}
\opt{B}{The mass can pass through the point (0, 0, 4).}
\opt{C}{If the position of the particle at t = 1s is (0, 2, 0), the x component of its momentum at t = 1 s is p$_\mathrm x$ = $\alpha\,\ell$n(2).}
\opt{D}{If the position of the particle at t = 1s is (0, 2, 0), the z component of its momentum at t = 1 s is p$_\mathrm z$ = 0.}
\end{qoptions}
\qn{2} A plane electromagnetic wave of the radiation field from a distant wire antenna travels toward the plane of the paper (which is in the -z direction). At time t = 0 s, the wave strikes the paper at normal incidence. The magnetic field vector at point O in the figure points in the -y direction and has a magnitude of 4.0 $\times$ 10$^{-8}$ T at that moment. This magnitude is maximum. The frequency of the wave is 1.0 $\times$ 10$^{16}$ Hz.\par
\qfig{p2_phy_s1_q02_1}
\begin{qoptions}
\opt{A}{The associated electric field E at time t = 0 s at the point O is 12$\hat{\mathrm i}$ V/m.}
\opt{B}{The associated electric field E at time t = 2.5 $\times$ 10$^{-17}$ s at z = 0.75 $\times$ 10$^{-8}$ m is 6$\sqrt2\hat{\mathrm i}$ V/m.}
\opt{C}{The associated electric field E at time t = 5 $\times$ 10$^{-17}$ s at z = 1.5 $\times$ 10$^{-8}$ m is $-$6$\sqrt2\hat{\mathrm i}$ V/m.}
\opt{D}{The associated electric field E at time t = 7.5 $\times$ 10$^{-17}$ s at z = 2.25 $\times$ 10$^{-8}$ m is $-$ 12$\hat{\mathrm i}$ V/m.}
\end{qoptions}
\qn{3} To an ideal battery with EMF 1.3 V we connected a bridge circuit assembled from three identical voltmeters and two identical milliammeters (see figure). It is known that the reading of one of the milliammeters is three times the other. Resistance of voltmeters are more than milliammeters.\par
\qfig{p2_phy_s1_q03_1}
\begin{qoptions}
\opt{A}{Reading of Voltmeter A is 1.0 V.}
\opt{B}{Reading of Voltmeter B is approx 0.62 V.}
\opt{C}{Reading of Voltmeter C is 0.4 V.}
\opt{D}{Reading of Voltmeter C is approx 0.68 V.}
\end{qoptions}
\qn{4} Consider a specially designed vernier calliper in which Vernier divisions are not equi spaced and main scale divisions are equispaced. 1 main scale division is equal to 1mm. Vernier scale has 8 divisions. When zero of the vernier scale is matched with zero of main scale, these 8 divisions of vernier scale coincide with 6 division of main scale division as shown.\par
\qfig{p2_phy_s1_q04_1}
When nothing is put between the jaws, zero of the main scale lies between zero and 1 division of Vernier scale and 4$^{\mathrm{th}}$ Vernier scale division exactly coincide with one of the main scale division. To measure the diameter of a cylindrical wire it is put between the jaws. Zero of vernier scale lies between 4$^{\mathrm{th}}$\& 5$^{\mathrm{th}}$ of main scale division and 6$^{\mathrm{th}}$ division of Vernier scale division coincide with one of the main scale division. Diameter of the wire is:
\begin{qoptions}
\opt{A}{Zero error is 0.2mm}
\opt{B}{Zero error is $-$0.2mm}
\opt{C}{Diameter is 4.3 mm}
\opt{D}{Diameter is 4.1 mm}
\end{qoptions}
\sect{SECTION-II(i)}
\paraq{1}{2}
A charged particle enters a region in which there is a frictional force proportional to the particle's velocity $(-\mathrm k\vec{\mathrm v})$, and the particle stops 10 cm from its entry point. If the particle repeats its motion when a homogeneous magnetic field of magnitude B, perpendicular to the plane of its trajectory, is also present, then the particle comes to rest at a displacement of 6 cm from its entry point (see figure).\par
(velocity with which particle enters the region is to be taken same in all cases)\par
\qfig{p2_phy_s2_para1_2_1}
\qn{1} If the magnetic field were twice as large the particle would stop at a displacement of $\dfrac{30}{\sqrt{\mathrm N}}$ cm from entry point, where N is:
\qn{2} The magnitude of the magnetic field for which the particle would stop at a displacement of 5$\sqrt2$ cm from entry point is $\dfrac{\mathrm{nB}}4$, where n is:
\paraq{3}{4}
Consider an adiabatic chamber, divided in three equal parts as shown. Pistons are massless and can slide freely inside the chamber, without any leakage of gas. Area of cross section of each piston is A. Piston-1 is conducting and piston-2 is adiabatic. Piston 1 and 2 are connected by a rigid massless rod whose thermal coefficient of linear expansion and heat capacity is negligible. Each part contains same gas, each at temperature T$_0$, pressure P$_0$ and volume V$_0$. Adiabatic exponent of gas is $\gamma$ = 2. Now using an electrical arrangement gas in part B is heated slowly so that its temperature becomes 2T$_0$.\par
\qfig{p2_phy_s2_para3_4_1}
\qn{3} Final temperature of gas in chamber A is nT$_0$ where n is:
\qn{4} Final pressure of gas in chamber A is P = $\left(\dfrac{\mathrm N+\sqrt{17}}8\right)$P$_0$, where N is :-
\paraq{5}{6}
A small container (1) of height 2$\ell_0$ and cross-section area `A' is kept in a big container (2) of height 3$\ell_0$ and cross-section area 2A. The bottom of container (1) is welded with the bottom of container (2). A massless piston is put on the container (1). The container (2) is completely filled with water and initially the length of air column in container (1) is $\ell_0$. Now a small orifice is made at the bottom of the container (2). Assume the temperature of water and air remains constant in all the following processes.\par
\qfig{p2_phy_s2_para5_6_1}
\qn{5} When the height of water level in container (2) becomes 2$\ell_0$, the piston reaches the topmost portion in container (1). The atmospheric pressure is N$\rho_\mathrm w$g$\ell_0$, where N is: (density of water is $\rho_\mathrm w$)\par
\qfig{p2_phy_s2_q05_1}
\qn{6} Now the piston is held (fixed) at that position. The orifice is closed and another orifice is made at the lowest position of container (1). The height up to which the water will rise in the container (1) at equilibrium is $\left(\mathrm N-\sqrt{\mathrm M}\right)\ell_0$ then (N + M) is :\par
\qfig{p2_phy_s2_q06_1}
\sect{SECTION-II(ii)}
\qn{7} Four identical condensers are connected as shown in figure, and joined to a battery E. Switch II is open to start with, and switch I closed. Switch I is then opened and switch II is closed.\par
\qfig{p2_phy_s3_q07_1}
{\footnotesize\setlength{\tabcolsep}{1.5pt}\noindent\begin{tabular}{|c|L{5.0cm}|c|L{1.0cm}|}
\hline
 & \multicolumn{1}{>{\centering\arraybackslash}p{5.0cm}|}{List-I} & & \multicolumn{1}{c|}{List-II} \\ \hline
(P) & ratio of energy stored in condenser `1' before and after closing switch II & (1) & 1 : 1 \\ \hline
(Q) & ratio of energy stored in condenser `2' before and after closing switch II & (2) & 1 : 4 \\ \hline
(R) & ratio of charge on condenser `2' before and after closing switch II & (3) & 4 : 1 \\ \hline
(S) & ratio of potential difference across condenser `3' before and after closing switch II & (4) & 2 : 1 \\ \hline
\end{tabular}}\par
\qn{8} Two infinite perpendicular planes were charged uniformly with surface densities of the charges being $\sigma$ and 2$\sigma$ respectively. A long hollow cylinder of radius R, whose axis coincides with the line of intersection of the planes, is given a surface charge $-$3$\sigma$, as shown in figure.\par
\qfig{p2_phy_s3_q08_1}
{\footnotesize\setlength{\tabcolsep}{1.5pt}\noindent\begin{tabular}{|c|L{4.2cm}|c|L{2.8cm}|}
\hline
 & \multicolumn{1}{>{\centering\arraybackslash}p{4.2cm}|}{List-I} & & \multicolumn{1}{>{\centering\arraybackslash}p{2.8cm}|}{List-II} \\ \hline
(P) & x-component of electric field at a point (x,y,z) inside cylinder & (1) & $\dfrac{\sigma}{2\varepsilon_0}\left[1-\dfrac{6\mathrm{Rx}}{(\mathrm x^2+\mathrm z^2)}\right]$ \\ \hline
(Q) & z-component of electric field at a point (x,y,z) inside cylinder & (2) & $\dfrac{\sigma}{2\varepsilon_0}$ \\ \hline
(R) & x-component of electric field at a point (x,y,z) outside cylinder & (3) & $\dfrac{\sigma}{2\varepsilon_0}\left[2-\dfrac{6\mathrm{Rz}}{(\mathrm x^2+\mathrm z^2)}\right]$ \\ \hline
(S) & z-component of electric field at a point (x,y,z) outside cylinder & (4) & $\dfrac{\sigma}{\varepsilon_0}$ \\ \hline
\end{tabular}}\par
\sect{SECTION-III}
\qn{1} Very thin non-conducting rod of length L is uniformly charged along the length with total charge of Q. A small conducting ring of radius R is made of very thin wire, its center coincides with one of the ends of the rod, and the plane of the ring is perpendicular to the rod. Ring has uniform charge q. With what net force F does the rod act on the ring? If the tension in ring due to charge on rod is T, find F/T. Round off to nearest integer. (Take L = $\sqrt3$R)
\qn{2} In an interference experiment as shown in the figure, the source plane and screen are separated by a distance 1m. At a certain position of source, fringe width is $\dfrac14$mm and by moving the source away from mirror along the line AB by 0.6 mm, the fringe width changed to $\dfrac16$mm, then wavelength of light used is (N $\times$ 1000) \AA, where N is:\par
\qfig{p2_phy_s4_q02_1}
\qn{3} A narrow beam of light is incident on a 30$^\circ-$60$^\circ-$90$^\circ$ prism kept in air as shown figure. The refractive index of the prism is $\mu$ = 2.1. The total clockwise deviation of the beam is $\dfrac{2\pi}{\mathrm N}$ radians, where N ($\neq$ 2) is:\par
\qfig{p2_phy_s4_q03_1}
\qn{4} A snap shot, of a vibrating string at t = 0, is shown in the figure. A particle at point P is observed moving up with velocity $\left(20\sqrt2\pi\right)$ cm/s and the tangent at the point P makes an angle $\theta=\tan^{-1}\left(2\sqrt2\pi\right)$ with the positive x-axis.\par
The equation of the wave is y = (4cm) sin $\left[\left(2\mathrm N\pi\,\mathrm s^{-1}\right)\mathrm t-\left(\dfrac{\mathrm M\pi}2\mathrm{cm}^{-1}\right)\mathrm x+\dfrac{\mathrm L\pi}4\right]$\par
(L $\to$ smallest positive integer). Find (N + M + L).\par
\qfig{p2_phy_s4_q04_1}
\qn{5} At low temperatures, the molar heat capacity of solids at constant pressure depends on the absolute temperature T according to the law C$_\mathrm T$ = C$_{\mathrm T_0}\left(\dfrac{\mathrm T}{\mathrm T_0}\right)^3$. One mole (40 g) of solid argon at temperature 9 K is brought into thermal contact with 600 gm of solid argon at a temperature of 1 K in an insulated chamber. What will be the temperature (in K) of the system when thermal equilibrium is established? The heat capacity of chamber is negligible. Round off to the nearest integer if necessary
\qn{6} In the photoelectric experiment, if we use a monochromatic light, the I-V curve is as shown. The work function of the metal is 2eV (Assume percentage efficiency of photo emission = 10$^{-3}$ \%, i.e. number of photoelectrons emitted are (10$^{-3}$) \% of number of photons incident on metal.) The power of light used in Watts is:\par
\qfig{p2_phy_s4_q06_1}
\raggedcolumns
\papersec{2}
\partsec{2}{CHEMISTRY}
\sect{SECTION-I}
\qn{1} In the following chemical reaction sequence which is/are correct.\par
\qfig{p2_chem_s1_q01_1}
\begin{qoptions}
\opt{A}{`P' gives yellow ppt. with I$_2$/OH$^-$.}
\opt{B}{`P' on treatment with NaBH$_4$ / EtOH will form a diol.}
\opt{C}{If 1 mole of `P' is treated with Sodium metal then 22.4 L of H$_2$(g) is released at 0$^\circ$C, 1 atm.}
\opt{D}{`P' is an aldehyde.}
\end{qoptions}
\qn{2} During the preparation of K$_2$Cr$_2$O$_7$, Chromite ore is fused with Na$_2$CO$_3$ in free excess of air. Then which is are correct?
\begin{qoptions}
\opt{A}{Chromite ore has iron in +2 oxidation state.}
\opt{B}{Oxidation number of Cr in fusion product is +6.}
\opt{C}{In the fusion product all `Cr$-$O' bond length are equal.}
\opt{D}{In the product Cr has no unpaired electron but still it is coloured.}
\end{qoptions}
\qn{3} \qfig{p2_chem_s1_q03_1}
Which is/are correct about above reaction sequence?
\begin{qoptions}
\opt{A}{`C' is a 2$^\circ$ alcohol.}
\opt{B}{`E' can't decolorize Br$_2$ + H$_2$O.}
\opt{C}{`C' to `D' conversion is nucleophilic substitution reaction.}
\opt{D}{Formation of `B' from `A' is dehydration reaction.}
\end{qoptions}
\qn{4} $\left(\Delta_\mathrm f^0\mathrm G\right)_{\mathrm{Cr}_2\mathrm O_{3(\mathrm s)}}$ = $-$527 kJ mol$^{-1}$\par
$\left(\Delta_\mathrm f^0\mathrm G\right)_{\mathrm{Al}_2\mathrm O_{3(\mathrm s)}}$ = $-$827 kJ mol$^{-1}$\par
$\left(\Delta_\mathrm f\mathrm H^0\right)_{\mathrm{Cr}_2\mathrm O_{3(\mathrm s)}}$ = $-$1100 kJ\par
$\left(\Delta_\mathrm f\mathrm H^0\right)_{\mathrm{Al}_2\mathrm O_{3(\mathrm s)}}$ = $-$1600 kJ\par
Which is/are correct?
\begin{qoptions}
\opt{A}{Cr$_2$O$_3$(s) is reduced by Al(s) in given standard condition.}
\opt{B}{Al$_2$O$_3$(s) is reduced by Cr(s) in given standard condition.}
\opt{C}{$\Delta_\mathrm r\mathrm H^\circ$ for the reduction of Cr$_2$O$_3$ by Al is $-$500 kJ.}
\opt{D}{Reduction of Cr$_2$O$_3$ by Al is spontaneous at all temperature.}
\end{qoptions}
\sect{SECTION-II(i)}
\paraq{1}{2}
Sulphate salt of an element `A' of 3-d series is (X). On heating `X' produces mixture of two gases. Both the gaseous product contains S$-$atom. Addition of few drops of HNO$_3$ to the solution of `X', followed by addition of KSCN produces blood red colour. Red colour disappears upon addition of excess of sodium oxalate, due to the formation of (Y). Then answer the following questions.
\qn{1} How many statement is/are true?\par
(i) Element `A' on treatment with dilute HCl forms salt in +3 oxidation state and releases H$_2$(g).\par
(ii) Decomposition of Sulphate salt (X) produces an oxide in +3 oxidation state of element.\par
(iii) Gases are SO$_2$ and SO$_3$.\par
(iv) Element is diamagnetic.\par
(v) `Y' is paramagnetic complex.
\qn{2} Find P + Q + R for compound (Y)\par
(P) Number of unpaired electron(s)\par
(Q) Number of stereoisomers\par
(R) Number of oxygen atoms
\paraq{3}{4}
The vapour pressure of two miscible liquids A and B are 300 and 500 mm Hg respectively. In a flask, 10 mole of A is mixed with 12 mole of B. However, as soon as `B' is added, `A' starts polymerizing into a completely insoluble solid. The polymerisation follows first order kinetics. After 100 min, 1 mole of a non$-$volatile solute is dissolved, which arrests the polymerisation completely. The final vapour pressure of the solution is 400 mm of Hg.\par
Assume negligible volume change on mixing and polymerisation and ideal behaviour for the final solution. [log2 = 0.3, log$_\mathrm e$10 = 2.3]
\qn{3} If rate constant for the first order is `K' min$^{-1}$, then find the value of $\left(\dfrac{10^4\mathrm K}{2.3}\right)$ is:
\qn{4} If initially mixed A and B has non$-$ideal behaviour with negative deviation from Raoult's law and there is no polymerisation, then how many of the following are true:\par
(i) Attraction between A and B is more than individual attraction of A.....A and B.....B.\par
(ii) Potential energy of the system decreases on mixing.\par
(iii) $\Delta$S$_{\mathrm{System}}$ $>$ 0\par
(iv) $\Delta$S$_{\mathrm{Surrounding}}$ $<$ 0\par
(v) $\Delta$G$_{\mathrm{Mixing}}$ $<$ 0
\paraq{5}{6}
Amino acids are connected by peptide bonds. A decapeptide (Mol. wt. 796) on complete hydrolysis gives glycine (Mol. wt. 75), alanine and phenylalanine. Glycine contributes 39.14\% to the total weight of the hydrolysed products. Then answer the following questions.
\qn{5} Find the total number of Glycine units present in the decapeptide.
\qn{6} This decapeptide consists of an octapeptide part with the following properties -\par
(i) All the glycine units are used in this octapeptide part \& it is single specific sequence\par
(ii) This octapeptide chain is fixed with an ordered linkage\par
(iii) The remaining two amino acids apart from this octapeptide are alanine and phenylalanine.\par
Then how many such decapeptides are possible
\sect{SECTION-II(ii)}
\qn{7} Match the Lists.\par
{\footnotesize\setlength{\tabcolsep}{1.5pt}\noindent\begin{tabular}{|c|c|c|c|}
\hline
 & \multicolumn{1}{c|}{List-I} & & \multicolumn{1}{c|}{List-II} \\ \hline
(P) & \qfigi{p2_chem_s3_q07_1} & (1) & \qfigi{p2_chem_s3_q07_5} \\ \hline
(Q) & \qfigi{p2_chem_s3_q07_2} & (2) & \qfigi{p2_chem_s3_q07_6} \\ \hline
(R) & \qfigi{p2_chem_s3_q07_3} & (3) & \qfigi{p2_chem_s3_q07_7} \\ \hline
(S) & \qfigi{p2_chem_s3_q07_4} & (4) & \qfigi{p2_chem_s3_q07_8} \\ \hline
\end{tabular}}\par
\qn{8} Match the Lists.\par
{\footnotesize\setlength{\tabcolsep}{1.5pt}\noindent\begin{tabular}{|c|L{3.8cm}|c|L{2.8cm}|}
\hline
 & \multicolumn{1}{>{\centering\arraybackslash}p{3.8cm}|}{List-I} & & \multicolumn{1}{>{\centering\arraybackslash}p{2.8cm}|}{List-II (pH)} \\ \hline
(P) & 10 ml of 0.1 M Na$_2$HPO$_4$ + 10 ml of 0.1 M HCl & (1) & pH = pKa$_1$ + log $\dfrac{[\mathrm{salt}]}{[\mathrm{acid}]}$ \\ \hline
(Q) & 10 ml of 0.1 M Na$_3$PO$_4$ + 15 ml 0.2 M HCl & (2) & pH = pKa$_2$ + log $\dfrac{[\mathrm{salt}]}{[\mathrm{acid}]}$ \\ \hline
(R) & 10 ml of 0.1 M H$_3$PO$_4$ + 15 ml 0.1 M NaOH & (3) & H$_3$PO$_4$ is the product \\ \hline
(S) & 15 ml of 0.1 M Na$_3$PO$_4$ + 40 ml of 0.1 M HCl & (4) & pH = $\dfrac12$(pKa$_2$ + pKa$_1$) \\ \hline
\end{tabular}}\par
\sect{SECTION-III}
\qn{1} To form a complete monolayer of acetic acid on 1g of charcoal, 100 mL of 0.5 M acetic acid was used. Some of the acetic acid remained unabsorbed. To neutralize the unabsorbed acetic acid, 40 mL of 1 M NaOH solution was required. If each molecule of acetic acid occupies \textbf{P} $\times$ 10$^{-21}$ m$^2$ surface area on charcoal, the value of $\left(\dfrac{\mathrm P}5\right)$ is \underline{\hspace{1.2cm}}.\par
[\textbf{Use given data:} Surface area of charcoal = 1.5 $\times$ 10$^2$ m$^2$g$^{-1}$; Avogadro's number (N$_\mathrm A$) = 6.0 $\times$ 10$^{23}$ mol$^{-1}$]
\qn{2} Given:\par
$\lambda^0_{\mathrm H^+}$ = 350 Scm$^2$mol$^{-1}$\par
$\lambda^0_{\mathrm A^-}$ = 50 Scm$^2$mol$^{-1}$\par
Molar conductivity of 0.038 M `HA' solution is 20 Scm$^2$mol$^{-1}$ at 298 K. What is the pH of 0.01 M NaA salt solution. \textbf{(Express your answer to nearest integer)}.
\qn{3} An organic compound (A) having molecular formula C$_5$H$_{10}$O responds +ve to Baeyers reagent as well as Lucas reagent. It shows geometrical isomerism as well as optical isomerism. (A) on heating with MnO$_2$ followed by treatment with Cl$_2$/NaOH forms (B) (organic product containing one carbon less than (A)). (B) when heated with NaOH/CaO forms organic compound (C). What is number of H$-$atoms in (C).
\qn{4} Find the number of combination(s) of reagents in which pink coloured solution or precipitate or product(s) is/are formed.\par
(i) MnCl$_2$ + H$_2$S(g) + HCl\qquad (ii) Mn$^{2+}$ + S$_2$O$_8^{2-}$\par
(iii) Mn$^{2+}$ + NaBiO$_3$ + HNO$_3$\qquad (iv) Mn$^{2+}$ + PbO$_2$ + HNO$_3$\par
(v) K$_2$MnO$_4$ + Cl$_2$\qquad (vi) ZnCl$_2$ + H$_2$S + NH$_4$OH(aq.)
\qn{5} Suppose de$-$Broglie equation is valid in Bohr model. In H$-$like sample, all electrons are in that orbit which has path length of the orbit equal to 3 times of de-Broglie wavelength of electrons in that orbit. This orbit of H$-$like sample has total energy equal to the ground state energy of H$-$atom. Now we have same H$-$like sample having all electrons in the ground state. If photons of energy 126.566 eV are given to this ground state sample and de-Broglie wavelength of the emitted electron is (P $\times$ 10$^{-10}$) metres then find nearest integer value of (P) ?
\qn{6} Consider the following reaction sequence.\par
\qfig{p2_chem_s4_q06_1}\par
Find the sum of number of stereoisomers of complexes P, Q and R.
\raggedcolumns
\papersec{2}
\partsec{3}{MATHEMATICS}
\sect{SECTION-I}
\qn{1} Which of the following is/are true.
\begin{qoptions}
\opt{A}{In any $\Delta$ABC, min value of\par{\small$\displaystyle\sum\dfrac{\sqrt{\sin A}}{\sqrt{\sin B}+\sqrt{\sin C}-\sqrt{\sin A}}$}\par is 3.}
\opt{B}{If $\sec A\tan B+\tan A\sec B=30$ then $\sec A\sec B+\tan A\tan B=\sqrt{901}$.}
\opt{C}{If $\sec A\tan B+\tan A\sec B=30$ then $\sec A\sec B+\tan A\tan B=\sqrt{899}$.}
\opt{D}{Three lines $y-z-1=0=x$, $z+x+1=0=y$, $x-z-1=0=y$ intersect the xy plane at U, M, R. Then the orthocenter of triangle UMR is $(0,1,0)$.}
\end{qoptions}
\qn{2} A is a symmetric matrix and B is a skew-symmetric matrix such that AB = BA, then
\begin{qoptions}
\opt{A}{$(A-B)^{-1}(A+B)$ is orthogonal matrix when A $-$ B is non-singular.}
\opt{B}{$(A+B)^{-1}(A-B)$ is orthogonal matrix when A + B is non-singular.}
\opt{C}{$\det\bigl((A-B)^{-1}(A+B)\bigr)=1$ and $\det\bigl((A+B)^{-1}(A-B)\bigr)=-1$}
\opt{D}{$\det\bigl((A-B)^{-1}(A+B)\bigr)=-1$ and $\det\bigl((A+B)^{-1}(A-B)\bigr)=1$}
\end{qoptions}
\qn{3} Let $y=f(x)$ be a differentiable curve passing through M(1, $-1$) such that the chord joining any two points A$(x_1,f(x_1))$ and B$(x_2,f(x_2))$ on the curve always meets the y-axis at N$(0,2x_1x_2)$. Then
\begin{qoptions}
\opt{A}{$\displaystyle\int_0^{1/2}f(x)\,dx=\dfrac1{12}$}
\opt{B}{$\displaystyle\int_0^{1/2}f(x)\,dx=\dfrac1{24}$}
\opt{C}{Area bounded by $f(x)$ and x-axis is $\dfrac{10}3$}
\opt{D}{Area bounded by $f(x)$ and x-axis is $\dfrac83$}
\end{qoptions}
\qn{4} Let a, b, c, d $\in$ R such that $\dfrac{a+2c}{b+3d}+\dfrac43=0$, then the equation $ax^3+bx^2+cx+d=0$ has
\begin{qoptions}
\opt{A}{at least one root in $(-1,0)$}
\opt{B}{at least one root in $(-1,1)$}
\opt{C}{at least one root in $(0,1)$}
\opt{D}{no root in $(-1,1)$}
\end{qoptions}
\sect{SECTION-II (i)}
\paraq{1}{2}
Let $f(x)=\displaystyle\lim_{n\to\infty}\Bigl(\cos\sqrt{\tfrac xn}\Bigr)^n$, $g(x)=\displaystyle\lim_{n\to\infty}\bigl(1-x+x\sqrt[n]{e}\bigr)^n$ and $h(x)=\tan^{-1}\bigl(g^{-1}(f^{-1}(x))\bigr)$\par
\qn{1} $4\left|\displaystyle\lim_{x\to0^+}\dfrac{\ell n(f(x))}{\ell n(g(x))}\right|$ is equal to :
\qn{2} Range of $y=h(x)$ is $(a,b)$, then $a+b$ is equal to :
\paraq{3}{4}
Let $P_n(x)=P_{n-1}(x-n)$, $n\in N\cup\{0\}$ and $P_0(x)=x^{90}-x^{89}+x^{88}-x^{87}+\ldots+1$.\par
\qn{3} If $P_{20}(x)=P_0(x-\lambda)$, then $\lambda$ equals :
\qn{4} If constant term in $P_{10}(x)$ is $\dfrac1p\bigl((55)^m-1\bigr)$, p and m $\in$ N, then p + m is equal to :
\paraq{5}{6}
Let $I(m)=\displaystyle\int_0^{\pi}\ell n\bigl(1-2m\cos x+m^2\bigr)dx$.\par
\qn{5} $\dfrac{I(81)}{I(3)}$ is :
\qn{6} $I(3)-I(-3)$ is :
\sect{SECTION-II (ii)}
\qn{7} Consider a triangle with vertices A$(z_1)\equiv$ A$\bigl(\tfrac2{\sqrt3}e^{i\pi/2}\bigr)$, B$(z_2)\equiv$ B$\bigl(\tfrac2{\sqrt3}e^{-i\pi/6}\bigr)$ and C$(z_3)\equiv$ C$\bigl(\tfrac2{\sqrt3}e^{-i5\pi/6}\bigr)$. A circle is inscribed in triangle ABC touching AB, BC and CA at D$(z_4)$, E$(z_5)$ and F$(z_6)$ respectively. Consider a point P(z) on the incircle of $\Delta$ABC.\par
{\footnotesize\setlength{\tabcolsep}{1.5pt}\noindent\begin{tabular}{|c|L{4.2cm}|c|L{1.6cm}|}\hline & \multicolumn{1}{>{\centering\arraybackslash}p{4.2cm}|}{List-I} & & \multicolumn{1}{>{\centering\arraybackslash}p{1.6cm}|}{List-II} \\ \hline
(P) & $\bigl|2\,\mathrm{Re}(z_4\bar z_5+z_5\bar z_6+z_6\bar z_4)\bigr|$ is & (1) & 4 \\ \hline
(Q) & $\dfrac{\text{Area }\Delta ABC}{\text{Area }\Delta DEF}$ is equal to & (2) & 1 \\ \hline
(R) & $AP^2+BP^2+CP^2$ is equal to & (3) & 5 \\ \hline
(S) & $DP^2+EP^2+FP^2$ is equal to & (4) & 2 \\ \hline
\end{tabular}}\par
\qn{8} Match the following List-I with List-II. Let $A=\begin{bmatrix}-3&-1&2\\3&1&-1\\4&2&5\end{bmatrix}$, $A\begin{bmatrix}x_1\\y_1\\z_1\end{bmatrix}=\begin{bmatrix}1\\0\\0\end{bmatrix}$, $A\begin{bmatrix}x_2\\y_2\\z_2\end{bmatrix}=\begin{bmatrix}0\\1\\0\end{bmatrix}$, $A\begin{bmatrix}x_3\\y_3\\z_3\end{bmatrix}=\begin{bmatrix}0\\0\\1\end{bmatrix}$ and $B=\begin{bmatrix}x_1&x_2&x_3\\y_1&y_2&y_3\\z_1&z_2&z_3\end{bmatrix}$.\par
{\footnotesize\setlength{\tabcolsep}{1.5pt}\noindent\begin{tabular}{|c|L{4.2cm}|c|L{1.6cm}|}\hline & \multicolumn{1}{>{\centering\arraybackslash}p{4.2cm}|}{List-I} & & \multicolumn{1}{>{\centering\arraybackslash}p{1.6cm}|}{List-II} \\ \hline
(P) & Trace B is & (1) & 256 \\ \hline
(Q) & 4 det(adj B) is & (2) & $-10$ \\ \hline
(R) & Det(adj(adj(2B))) is & (3) & 1 \\ \hline
(S) & Sum of element of B is & (4) & $-8$ \\ \hline
\end{tabular}}\par
\sect{SECTION-III}
\qn{1} E$_1$ and E$_2$ are ellipses and the area of E$_2$ is one-third the area of the quadrilateral formed by tangents at the ends of the latus rectum of ellipse E$_3$: $5x^2+9y^2=45$. E$_1$, E$_2$, E$_3$ are similar and concentric. E$_1$ is inscribed in E$_2$ so that both touch each other and the minor axis of E$_2$ coincides with the major axis of E$_1$. If the length of the major axis of E$_1$ is equal to the length of the minor axis of E$_2$, then the area of E$_2$ outside E$_1$ is equal to :
\qn{2} Let $f(x)$ be a fifth order polynomial and every root of $f(x)=0$ is real and distinct. The number of real roots of $f''(x)f(x)-(f'(x))^2=0$ is :
\qn{3} In an organization, the number of women is $\mu$ times that of men. If $\alpha$ things are distributed among them, the probability that number of things received by men is odd is $\left(\dfrac12-\left(\dfrac12\right)^{\alpha+1}\right)$. Then find the value of $\mu$ :
\qn{4} Let $\vec a,\vec b,\vec c$ be vectors of magnitudes 1, 2, 3 respectively and $\bigl|[\vec a\ \vec b\ \vec c]\bigr|=6$. If $\vec r$ is a unit vector coplanar with $\vec b$ and $\vec c$ and $\vec r\cdot\vec b=1$, then the value of $\bigl|(\vec a\times\vec c)\times\vec r\bigr|^2+\bigl|(\vec a\times\vec c)\cdot\vec r\bigr|^2$ is :
\qn{5} On the side AC of an acute triangle ABC a point D is taken such that AD = 1, DC = 2 and BD is an altitude of the triangle. A circle of radius 2 passes through the points A and D and touches at the point D the circle circumscribed about the triangle BDC. The area of $\Delta$ABC is $p\sqrt q$ ($p$ is prime and $q$ is natural number). Then the value of $\dfrac{q-p}2$ is :
\qn{6} Let $S=\displaystyle\sum_{r=1}^{99}\dfrac{2r^2-98r+1}{(100-r)\cdot{}^{100}C_r}$. Then the value of $\dfrac{[S]}{11}$ is (where [.] is the greatest integer function) :
\end{multicols}
\end{document}
